\section{Gramatika jazyka}
\label{sec:gramatika}

Tato sekce popisuje implementovanou syntaxi jazyka \Roth{} ve zjednodušené EBNF.
Jazyk je tokenizován na základě bílých znaků a několika speciálních znaků
(\texttt{:}, \texttt{;}, \texttt{(}, \texttt{)}, \texttt{"}).

\subsection{Lexikální pravidla}

\begin{description}
	\item[Whitespace] mezery, tabulátory a konce řádků oddělují tokeny.
	\item[Comment] \texttt{(} \emph{libovolné znaky bez \texttt{)}} \texttt{)}. Komentáře jsou ignorovány parserem.
	\item[Number] dekadický literál \texttt{i32}. Pokud převod selže (např. přetečení), token je brán jako slovo.
	\item[Word] posloupnost znaků do prvního bílého znaku, \texttt{(}, \texttt{)}, \texttt{"}. Slova jsou case-insensitive (lexer je normalizuje na velká písmena).
	\item[String] buď \texttt{"..."} (s podporou escape sekvencí \texttt{\\n}, \texttt{\\t}, \texttt{\\r}, \texttt{\\"}, \texttt{\\\\}), nebo \texttt{S" ..."} (bez escape sekvencí; volitelná mezera po \texttt{S"}).
\end{description}

\subsection{Syntaktická pravidla (EBNF)}

\begin{lstlisting}[language=Forth, caption={EBNF gramatika}]
program        ::= { statement } ;

statement      ::= number
                | string
                | word
                | definition
                | variable_decl
                | comment ;

definition     ::= ':' word { stmt_in_def } ';' ;
stmt_in_def    ::= number | string | word | variable_decl | comment ;
                // vnorene definice nejsou dovoleny

variable_decl  ::= 'VARIABLE' word ;

number         ::= [ '-' ] digit { digit } ;

string         ::= '"' { char | escape } '"'
                | 'S"' [ ' ' ] { char } '"' ;

comment        ::= '(' { char } ')' ;

word           ::= token ;
\end{lstlisting}

Při kompilaci ze souboru je před lexikální analýzou proveden preprocesing
direktivy \texttt{INCLUDE}, která rekurzivně vkládá obsah souborů do zdrojového
textu. Název souboru může být uveden v uvozovkách nebo bez nich a cesta je
vyhodnocena relativně k aktuálnímu souboru (a případně k pracovnímu adresáři).
